% from ega2-1-fr.tex
\oldpage[II]{11}
$h:Z\to S$ est égal à $g\circ f\circ p$ et à $g\circ f'\circ p'$. On
définit un homomorphisme canonique
\[
  \mathcal{A}(\mathcal{F})\otimes_{\mathcal{A}(Y)}\mathcal{A}(\mathcal{F}')
  \longrightarrow \mathcal{A}(\mathcal{F}\otimes_Y\mathcal{F}')
\]
de la façon suivante : pour tout ouvert $U\subset S$, on a des homomorphismes
canoniques
\[
  \begin{aligned}
  \Gamma(f^{-1}(g^{-1}(U)),\mathcal{F})
  &\longrightarrow \Gamma(h^{-1}(U),p^*(\mathcal{F}))\quad\text{et}\quad
  \Gamma(f^{\prime-1}(g^{-1}(U)),\mathcal{F}')
  \longrightarrow \Gamma(h^{-1}(U),p^{\prime*}(\mathcal{F}')).
  \end{aligned}
\]
(\hyperref[0.4.4.3-fr]{0, 4.4.3}), d'où on déduit un homomorphisme canonique
\[
  \begin{aligned}
  &\Gamma(f^{-1}(g^{-1}(U)),\mathcal{F})
    \otimes_{\Gamma(g^{-1}(U),\mathcal{O}_Y)}
    \Gamma(f^{\prime-1}(g^{-1}(U)),\mathcal{F}')\longrightarrow\\
  &\hspace{17em}
    \Gamma(h^{-1}(U),p^*(\mathcal{F}))
    \otimes_{\Gamma(h^{-1}(U),\mathcal{O}_Z)}
    \Gamma(h^{-1}(U),p^{\prime*}(\mathcal{F}')).
  \end{aligned}
\]

Pour voir qu'on a bien là un isomorphisme de $\mathcal{A}(Z)$-Modules, on
peut se borner au cas où $S$ (et par suite $X$, $X'$, $Y$,
$X\times_Y X'$) sont des schémas affines, et (avec les notations de
(\hyperref[II.1.4.6-fr]{1.4.6})), $\mathcal{F}=\widetilde{M}$,
$\mathcal{F}'=\widetilde{M'}$, $M$ (resp. $M'$) étant un $A$-module (resp. un
$A'$-module). Alors $\mathcal{F}\otimes_Y\mathcal{F}'$ s'identifie au
faisceau sur $X\times_Y X'$ associé au $(A\otimes_B A')$-module
$M\otimes_B M'$ (\hyperref[I.9.1.3-fr]{I, 9.1.3}) et le corollaire résulte de
l'identification canonique des $\mathcal{O}_S$-Modules
$\widetilde{M\otimes_B M'}$ et
$\widetilde{M}\otimes_{\widetilde{B}}\widetilde{M'}$ (où $M$, $M'$ et $B$
sont considérés comme des $C$-modules)
(\hyperref[I.1.3.13-fr]{I, 1.3.13} et
\hyperref[I.1.6.3-fr]{I, 1.6.3}).

Si on applique en particulier (\hyperref[II.1.4.7-fr]{1.4.7}) au cas où
$X=Y$ et $\mathcal{F}'=\mathcal{O}_{X'}$, on voit que le
$\mathcal{A}'$-Module $\mathcal{A}(f^{\prime*}(\mathcal{F}))$ s'identifie à
$\mathcal{A}(\mathcal{F})\otimes_{\mathcal{B}}\mathcal{A}'$.

\begin{env}[1.4.8]
\label{II.1.4.8-fr}
En particulier, lorsque $X=X'=Y$ ($X$ étant affine au-dessus de $S$), on voit
que si $\mathcal{F}$, $\mathcal{G}$ sont deux $\mathcal{O}_X$-Modules
quasi-cohérents, on a
\[
  \mathcal{A}(\mathcal{F}\otimes_{\mathcal{O}_X}\mathcal{G})
  =\mathcal{A}(\mathcal{F})\otimes_{\mathcal{A}(X)}\mathcal{A}(\mathcal{G})
  \tag{1.4.8.1}\label{II.1.4.8.1-fr}
\]
à un isomorphisme canonique fonctoriel près. Si de plus $\mathcal{F}$ admet
une présentation finie, il résulte de
(\hyperref[I.1.6.3-fr]{I, 1.6.3} et
\hyperref[I.1.3.12-fr]{I, 1.3.12}) que
\[
  \mathcal{A}(\mathcal{H}om_X(\mathcal{F},\mathcal{G}))
  =\mathcal{H}om_{\mathcal{A}(X)}
    (\mathcal{A}(\mathcal{F}),\mathcal{A}(\mathcal{G}))
  \tag{1.4.8.2}\label{II.1.4.8.2-fr}
\]
à un isomorphisme canonique près.
\end{env}

\begin{remark}[1.4.9]
\label{II.1.4.9-fr}
Si $X$, $X'$ sont deux préschémas affines au-dessus de $S$, la somme
$X\sqcup X'$ est aussi affine au-dessus de $S$, la somme de deux schémas
affines étant un schéma affine.
\end{remark}

\begin{proposition}[1.4.10]
\label{II.1.4.10-fr}
Soient $S$ un préschéma, $\mathcal{B}$ une $\mathcal{O}_S$-Algèbre
quasi-cohérente, $X=\operatorname{Spec}(\mathcal{B})$. Pour tout faisceau
quasi-cohérent d'idéaux $\mathcal{J}$ de $\mathcal{B}$,
$\widetilde{\mathcal{J}}$ est un faisceau quasi-cohérent d'idéaux de
$\mathcal{O}_X$, et le sous-préschéma fermé $Y$ de $X$ défini par
$\widetilde{\mathcal{J}}$ est canoniquement isomorphe à
$\operatorname{Spec}(\mathcal{B}/\mathcal{J})$.
\end{proposition}

En effet, il résulte aussitôt de (\hyperref[I.4.2.3-fr]{I, 4.2.3}) que $Y$
est affine au-dessus de $S$~; en vertu de
(\hyperref[II.1.3.1-fr]{1.3.1}), on est donc ramené au cas où $S$ est affine,
et la proposition résulte alors aussitôt de
(\hyperref[I.4.1.2-fr]{I, 4.1.2}).

On peut encore exprimer le résultat de
(\hyperref[II.1.4.10-fr]{1.4.10}) en disant que si
$h:\mathcal{B}\to\mathcal{B}'$ est un homomorphisme \emph{surjectif} de
$\mathcal{O}_S$-Algèbres quasi-cohérentes,
${}^a h:\operatorname{Spec}(\mathcal{B}')\to
\operatorname{Spec}(\mathcal{B})$ est une \emph{immersion fermée}.

